\section{Felhasználói felület}

A felhasználói felület nem egyetlen összetett nézetben tartalmazza az alkalmazás összes funkcióját,
hanem a munkamenet fő lépéseihez igazodva külön képernyőkre bontja azokat.
Ez a felépítés segíti az áttekinthetőséget, mivel a felhasználó minden lépésben csak az adott
művelethez szükséges információkkal és vezérlőelemekkel találkozik.


\subsection{Konfigurációs képernyő}

A konfigurációs képernyő a munkamenet előkészítésére szolgál. Itt hozható létre a
Bluetooth-kapcsolat a Sensoria okoszoknival, amelynek állapotáról az alkalmazás visszajelzést
ad a felhasználónak. A sikeres csatlakozás után a felhasználó megadhatja a nevét,
amely a későbbiekben a mérés azonosítására szolgál.


A rehabilitációs mozgássorozat kétféleképpen állítható össze:
\begin{itemize}
    \item Egy előre meghatározott gyakorlatsor használata, amely egy specifikusan
    összeállított rehabilitációs feladatsort tartalmaz. Ebben az esetben a gyakorlatok sorrendje
    adott, a hozzájuk tartozó időtartamok azonban módosíthatók.
    \item Egyéni feladatsor összeállítása. Ez akkor hasznos, ha az előre definiált gyakorlatsor
    nem illeszkedik a felhasználó aktuális rehabilitációs szakaszához vagy állapotához.
    Ilyenkor a felhasználó egy legördülő listából választhatja ki a kívánt gyakorlatokat,
    megadhatja azok időtartamát, majd tetszőleges sorrendben hozzáadhatja őket a munkamenethez.
\end{itemize}


A konfigurációs képernyőről indítható el a mérés, valamint innen érhetők el a
korábban mentett munkamenetek is. A képernyő így a mérési folyamat kiindulópontjaként működik.


\subsection{Végrehajtási képernyő}

A végrehajtási képernyő a konfigurációs képernyőről indított mérés után jelenik meg.
A mérés megkezdése előtt a felhasználó egy felugró ablakban (dialog) választhatja ki a
kívánt vizuális visszacsatolási módot.
Lehetőség van csak a feladatot bemutató videó megjelenítésére,
csak az előlapi kamera képének megjelenítésére, illetve a kettő egyidejű,
egymás melletti megjelenítésére.


Minden gyakorlat előtt egy rövid, öt másodperces előkészítő szakasz jelenik meg.
Ebben az időszakban a felhasználó a gyakorlatot bemutató videót látja,
így még a mérés kezdete előtt megfigyelheti a helyes végrehajtást.
A tényleges adatgyűjtés és a gyakorlat végrehajtása csak ezt követően indul el.


A mérés a munkamenet közben bármikor megszakítható. Ebben az esetben az alkalmazás nem menti el
a munkamenetet, így nem keletkezik hiányos vagy félbeszakított adatokon alapuló elemzés.
Ez fontos szempont az eredmények értelmezhetősége és a tárolt adatok következetessége miatt.


A gyakorlatok végrehajtása közben a felület a kiválasztott megjelenítési módnak megfelelő
videós tartalmat jeleníti meg. A videó alatt a statikus szöveges instrukciók szerepelnek,
illetve a stabilitási adatok jelennek meg, amelyek az aktuális végrehajtásról adnak közvetlen
visszajelzést a felhasználónak.


\subsection{Eredményképernyő}

A mérés lezárása után az alkalmazás automatikusan az eredményképernyőre navigál.
Ez a felület a munkamenethez tartozó összesített információkat jeleníti meg strukturált formában.
A cél az, hogy a rögzített adatok ne nyers szenzorértékekként jelenjenek meg, hanem olyan
áttekinthető formában, amely segíti a mérés értelmezését.


Az eredményképernyőn megjelennek a munkamenet azonosításához szükséges alapadatok,
például a felhasználó neve és a mérés időpontja.
Emellett láthatók a végrehajtott gyakorlatokhoz kapcsolódó eredmények,
valamint a szenzoradatokból számított mutatók. Ezek az információk a munkamenet későbbi
áttekintését és összehasonlítását támogatják.

\subsection{Mentett munkamenet képernyő}

A mentett munkamenetek képernyője a korábban rögzített mérések elérésére szolgál.
A felület listában jeleníti meg az elmentett munkameneteket, amelyek közül a felhasználó
kiválaszthatja a megtekinteni kívánt elemet. A kiválasztott munkamenet betöltése után az
alkalmazás megjeleníti az ahhoz tartozó eredményképernyőt.


A listázás során  az egyes mérések egyértelműen azonosíthatók a felhasználó neve és a
méréshez tartozó időbélyeg segítségével.
Így több eltárolt munkamenet esetén is visszakereshető, hogy melyik mérés kihez és mikorhoz
kapcsolódik.


A nem kívánt vagy már nem szükséges munkamenetek manuálisan törölhetők a telefon tárhelyéről.
Ez lehetőséget ad arra, hogy a felhasználó vagy a fejlesztő szükség esetén eltávolítsa a már nem
releváns mérési adatokat.

